% Vietnamese translation for OI Public Library
% Source-File: IOI中国国家候选队论文/2004/黄源河.ppt
\documentclass[11pt]{article}
\usepackage{oipl-vn}

\title{Xây dựng và ứng dụng mô hình đồ thị\\{\large Ghi chú theo slide}}
\author{Huang Yuanhe}
\date{IOI 2004}

\begin{document}
\maketitle

\origtitle{Building and applying graph models}
\sourcepath{IOI 2004 / Huang Yuanhe.ppt}

\section{Mô hình hóa bằng đồ thị}

Các slide dùng một số bài để minh họa việc biến quan hệ ràng buộc thành đỉnh,
cạnh, luồng hoặc ghép cặp. Khi mô hình đúng, bài toán gốc thường trở thành
một bài đồ thị kinh điển; khi mô hình sai, thuật toán đồ thị mạnh cũng không
cứu được lời giải.

\section{Place the Robots}

Ví dụ đầu là \emph{Place the Robots}, với lưới \(N\times M\),
\(N,M\le 50\). Ý tưởng quen thuộc là chia các ô trống liên tiếp theo hàng và
theo cột thành các đoạn, rồi dựng đồ thị hai phía giữa đoạn hàng và đoạn cột
nếu chúng giao nhau tại một ô có thể đặt robot. Khi đó số robot tối đa là
kích thước ghép cực đại.

\section{Thuê thu ngân}

Bài ACM Tehran 2000 xuất hiện với các biến \(R_0,R_1,\ldots,R_{23}\), chỉ số
\(i=0,\ldots,23\) và ràng buộc \(W_i\), \(N\le 1000\). Đây là dạng bài chênh
lệch trên vòng thời gian 24 giờ: đặt biến tổng tiền tố cho số người được thuê
tới từng giờ, rồi chuyển yêu cầu phục vụ thành các cạnh bất đẳng thức. Sau đó
có thể dùng đường đi ngắn nhất hoặc kiểm tra khả thi hệ ràng buộc.

\section{Greedy Island}

Ví dụ cuối nhắc ngưỡng \(100000\) và biểu thức \(A+B+C\). Nội dung chính là
chọn mô hình đồ thị sao cho quyết định tham lam trên bài gốc được thay bằng
quan hệ đường đi hoặc kết nối rõ ràng trong đồ thị.

\section{Ghi nhớ}

Mô hình đồ thị tốt cần chỉ ra: đỉnh biểu diễn gì, cạnh biểu diễn quan hệ nào,
trọng số hoặc dung lượng mang ý nghĩa gì, và nghiệm đồ thị được đổi ngược về
nghiệm bài gốc ra sao.

\end{document}
